\begin{exercise}
From the text, Exercise 3.2.
\end{exercise}


\begin{exercise}
From the text, Exercise 3.11.
\end{exercise}

%\begin{exercise}
%A {\bf complete matching} of a graph with $2n$ verticies is a subgraph consisting of $n$ disjoint edges.  How many complete matching of this graph?\\
%\begin{figure}[H]
%\centering
%\includegraphics[scale=2]{3dot11} 
%\end{figure}
%\end{exercise}


%\begin{exercise}
%A tree with $n$ vertices is called {\bf graceful} if it's vertices can be labeled with the integers 1, 2 , ..., $n$ such that the absolute values of the difference of the labels of adjacent vertices are all different.  Draw and label two non-isomorphic graceful trees on 6 vertices.
%\end{exercise}

%\begin{exercise}
%A {\bf forest} is a simple graph where each component is a tree.  Find all the non-isomorphic forests with five vertices and two or more components.
%\end{exercise}

\begin{exercise} 
Find the definition for {\bf isomorphic} graphs.  What is the definition?  Give an example of two graphs that are isomorphic, and give an example of two graphs that are not isomorphic.
\end{exercise}


\begin{exercise} 
A simple graph $G$ is {\bf self-complimentary} if $G$ and $\bar{G}$ are isomorphic.  Find a self-complimentary graph having five vertices.  Find another self complimentary graph with more than 1 vertex.
\end{exercise}

%\begin{exercise}
%For the following degree sequences (list of the degrees of the vertices in a graph) determine if there exists a bipartite simple graph with these degrees.  If yes construct it, if not explain why it does not exist.  (a) (6, 6, 4, 4, 4, 4, 4, 4, 2, 2, 2); (b) (4, 4, 4, 4, 4, 3, 3, 2, 2); (c) (5, 5, 5, 5, 5, 5, 4, 4, 4)
%\end{exercise}


\begin{exercise}
A vertex $v$ in a connected graph $G$ is an {\bf articulation point} if the removal of $v$ and all edges incident to $v$ disconnects $G$.  Give an example of a graph with six vertices that has exactly two articulation points.  Give an example of a graph on six vertices that has no articulation points.
\end{exercise}

%\begin{proofp}
%Prove that if there are two different paths between vertices $x$ and $y$ in a simple graph, then there must be a cycle in the graph.
%\end{proofp}

%\begin{proofp}
%Prove that if a simple graph contains an odd length closed trail $v_0 \rightarrow v_1 \rightarrow v_2 \rightarrow \cdots \rightarrow v_k$ then the graph contains an odd length cycle.
%\end{proofp}


\begin{proofp}
Prove that if G is a simple graph with $p$ vertices, where each vertex has degree $\ge \frac{1}{2}(p-1)$, then $G$ must be connected.  Hint: Try a contradiction proof, by assuming there are at least two components of $G$.  How many vertices must each component have?
\end{proofp}

\begin{proofp}
Prove the statement: In any tree there is exactly one path from any vertex to any other vertex. (Hint: Try contradiction. If there are two paths try to show there is a cycle.)
\end{proofp}



\begin{proofp}
A {\bf forest} is a simple graph where each component is a tree.  Find a formula for the number of edges in a forest with $n$ vertices and $c$ components.  Prove your formula. 
\end{proofp}

\begin{proofp}
Prove the average degree in a tree is always less than 2. More specifically express this average as a function of the number of vertices in tree.
\end{proofp}



\begin{proofp}$\bigstar$
Let $v$ and $w$ be distinct vertices in the complete graph $K_n$ for $n \ge 2$.  Prove the number of paths from $v$ to $w$ is
$$(n-2)!\sum^{n-2}_{k=0}\frac{1}{k!}$$
\end{proofp}

\begin{proofp}$\bigstar$
Prove that a vertex $v$ in a connected graph $G$ is an articulation point if and only if there are vertices $w$ and $x$ in $G$ having the property that every path from $w$ to $x$ passes through $v$.
\end{proofp}



\begin{proofp}$\bigstar$
Prove that in any simple graph with more than one vertex, there must be two vertices with the same degree. (Hint: Try contradiction.)
\end{proofp}



\begin{proofp}$\bigstar$
Prove the number of vertices of degree 1 in an tree must be greater than or equal to the maximum degree in the tree.
\end{proofp}